\begin{theorem}[Finite stationary-phase duality]\label{mod:lamprecht-finiteduality}
Let $E$ be a nonarchimedean local field and let
$\psi:E\to\C^\times$ be a continuous additive character with exact
conductor $n=n_E(\psi)$.  Let $M,q,r\in\mathbf Z$ satisfy $r\le q$, and
let $\Gamma\in E^\times$ satisfy
\[
 \operatorname{ord}_E(\Gamma)=M+n.
\]
Put
\[
 A_{M,q,r}=\mathfrak p_E^{M-q}/\mathfrak p_E^{M-r},
 \qquad
 B_{q,r}=\mathfrak p_E^r/\mathfrak p_E^q.
\]
There is a representative-independent biadditive pairing
\begin{equation}\label{eq:blueprint-general-stationary-pairing}
 \langle\ ,\ \rangle_{\psi,\Gamma}:
 A_{M,q,r}\times B_{q,r}\longrightarrow\C^\times,
 \qquad
 \langle\bar c,\bar x\rangle_{\psi,\Gamma}
   =\psi(cx/\Gamma).
\end{equation}
More precisely, multiplication first gives a bilinear quotient map
\[
 A_{M,q,r}\times B_{q,r}
 \longrightarrow
 \mathfrak p_E^{M-q+r}/\mathfrak p_E^M,
 \qquad (\bar c,\bar x)\longmapsto\overline{cx},
\]
because
\[
 M-q+r\le M,
 \qquad M=(M-r)+r=(M-q)+q.
\]
The character $z\mapsto\psi(z/\Gamma)$ descends through the denominator
$\mathfrak p_E^M$, since division by $\Gamma$ carries that denominator to
$\mathfrak p_E^{-n}$, where $\psi$ is trivial.  Thus no representative of
either input quotient is selected by the construction.

The two annihilators occur at the exact stated depths:
\[
 \begin{aligned}
  \bigl(\forall x\in\mathfrak p_E^r,
     \ \psi(cx/\Gamma)=1\bigr)
    &\Longleftrightarrow c\in\mathfrak p_E^{M-r},\\
  \bigl(\forall c\in\mathfrak p_E^{M-q},
     \ \psi(cx/\Gamma)=1\bigr)
    &\Longleftrightarrow x\in\mathfrak p_E^q.
 \end{aligned}
\]
Equivalently, on quotient classes,
\[
 \begin{aligned}
  \bigl(\forall\bar x\in B_{q,r},
    \ \langle\bar c,\bar x\rangle_{\psi,\Gamma}=1\bigr)
    &\Longleftrightarrow\bar c=0,\\
  \bigl(\forall\bar c\in A_{M,q,r},
    \ \langle\bar c,\bar x\rangle_{\psi,\Gamma}=1\bigr)
    &\Longleftrightarrow\bar x=0.
 \end{aligned}
\]
The forward implications use an element $w$ of exact order $-n-1$ with
$\psi(w)\ne1$, supplied by exact conductor data.  In particular, they do
not assume the false pointwise converse
$\psi(z)=1\Rightarrow z\in\mathfrak p_E^{-n}$.

Both quotients are finite and have the same exact cardinality
\[
 |A_{M,q,r}|=|B_{q,r}|=q_E^{(q-r).\operatorname{toNat}}.
\]
After composing $\C^\times\hookrightarrow\C$, currying the pairing in
either variable therefore gives bijective additive homomorphisms
\[
 A_{M,q,r}\xrightarrow{\ \sim\ }
   \operatorname{AddChar}(B_{q,r},\C),
 \qquad
 B_{q,r}\xrightarrow{\ \sim\ }
   \operatorname{AddChar}(A_{M,q,r},\C).
\]
The Lean theorem \texttt{lamprechtPerfectPairing} is the conjunction of
these two bijectivity assertions, and the two resulting additive
equivalences are exported separately.

For Lamprecht's classical parameters $m=2d+\varepsilon$ and
$r=d+\varepsilon$, specialize only by taking $M=q=m$ and
$\Gamma=\gamma$.  Then
\[
 A_{m,m,d+\varepsilon}=\mathcal O_E/\mathfrak p_E^d,
 \qquad
 B_{m,d+\varepsilon}
   =\mathfrak p_E^{d+\varepsilon}/\mathfrak p_E^m,
\]
and \eqref{eq:blueprint-general-stationary-pairing} is precisely the
pairing in the manuscript's finite-duality lemma.
\LeanFile{LanglandsFirstMainLemma/Lamprecht/FiniteDuality.lean}
\lean{LanglandsFirstMainLemma.LamprechtCoefficientQuotient}
\lean{LanglandsFirstMainLemma.LamprechtVariableQuotient}
\lean{LanglandsFirstMainLemma.lamprechtAnnihilatorLeft}
\lean{LanglandsFirstMainLemma.lamprechtAnnihilatorRight}
\lean{LanglandsFirstMainLemma.lamprechtProductCharacter}
\lean{LanglandsFirstMainLemma.lamprechtPairing}
\lean{LanglandsFirstMainLemma.lamprechtPairing_mk_mk}
\lean{LanglandsFirstMainLemma.lamprechtPairingLeft}
\lean{LanglandsFirstMainLemma.lamprechtPairingRight}
\lean{LanglandsFirstMainLemma.lamprechtPairingLeft_eq_zero_iff}
\lean{LanglandsFirstMainLemma.lamprechtPairingRight_eq_zero_iff}
\lean{LanglandsFirstMainLemma.lamprechtPairing_left_nondegenerate}
\lean{LanglandsFirstMainLemma.lamprechtPairing_right_nondegenerate}
\lean{LanglandsFirstMainLemma.lamprechtCoefficientQuotient_card}
\lean{LanglandsFirstMainLemma.lamprechtVariableQuotient_card}
\lean{LanglandsFirstMainLemma.lamprechtPerfectPairing}
\lean{LanglandsFirstMainLemma.lamprechtPairingLeftEquiv}
\lean{LanglandsFirstMainLemma.lamprechtPairingRightEquiv}
\leanok
\SourceText{Lemma lem:finite-duality; equation eq:general-stationary-pairing.}
\uses{mod:localfield-finitequotients,mod:basic-characterconductors}
\FormalizationContract The denominator is supplied as an element of $E^\times$, exact
conductor data provide both triviality and the predecessor witness, and every pairing
input is a quotient class.  No pointwise-kernel characterization and no quotient
representative choice is assumed.
\end{theorem}
